\documentclass[aspectratio=169,11pt]{beamer}
\newcommand{\footerlabel}{DEMO101 Algorithms \textbullet{} Lecture 2}
% 共用样式 —— 虚构示例课件（DEMO），内容全部自编，不来自任何真实课程。
% 用法：每份 deck 先 \documentclass，再定义 \footerlabel，最后 % 共用样式 —— 虚构示例课件（DEMO），内容全部自编，不来自任何真实课程。
% 用法：每份 deck 先 \documentclass，再定义 \footerlabel，最后 % 共用样式 —— 虚构示例课件（DEMO），内容全部自编，不来自任何真实课程。
% 用法：每份 deck 先 \documentclass，再定义 \footerlabel，最后 \input{_style}。
\usepackage[T1]{fontenc}
\usepackage{lmodern}
\usepackage{amsmath,amssymb}
\usetheme{default}
\setbeamertemplate{navigation symbols}{}
\setbeamercolor{frametitle}{fg=blue!50!black}
\setbeamerfont{frametitle}{series=\bfseries,size=\Large}
\setbeamercolor{title}{fg=blue!50!black}
\setbeamercolor{footline}{fg=black!60}
\setbeamertemplate{footline}{%
  \leavevmode\hbox{%
    \begin{beamercolorbox}[wd=.85\paperwidth,ht=2.5ex,dp=1ex,left,leftskip=1em]{footline}\footnotesize\footerlabel\end{beamercolorbox}%
    \begin{beamercolorbox}[wd=.15\paperwidth,ht=2.5ex,dp=1ex,right,rightskip=1em]{footline}\footnotesize\insertframenumber\end{beamercolorbox}}%
  \vskip0pt}
\setbeamertemplate{itemize items}[circle]
。
\usepackage[T1]{fontenc}
\usepackage{lmodern}
\usepackage{amsmath,amssymb}
\usetheme{default}
\setbeamertemplate{navigation symbols}{}
\setbeamercolor{frametitle}{fg=blue!50!black}
\setbeamerfont{frametitle}{series=\bfseries,size=\Large}
\setbeamercolor{title}{fg=blue!50!black}
\setbeamercolor{footline}{fg=black!60}
\setbeamertemplate{footline}{%
  \leavevmode\hbox{%
    \begin{beamercolorbox}[wd=.85\paperwidth,ht=2.5ex,dp=1ex,left,leftskip=1em]{footline}\footnotesize\footerlabel\end{beamercolorbox}%
    \begin{beamercolorbox}[wd=.15\paperwidth,ht=2.5ex,dp=1ex,right,rightskip=1em]{footline}\footnotesize\insertframenumber\end{beamercolorbox}}%
  \vskip0pt}
\setbeamertemplate{itemize items}[circle]
。
\usepackage[T1]{fontenc}
\usepackage{lmodern}
\usepackage{amsmath,amssymb}
\usetheme{default}
\setbeamertemplate{navigation symbols}{}
\setbeamercolor{frametitle}{fg=blue!50!black}
\setbeamerfont{frametitle}{series=\bfseries,size=\Large}
\setbeamercolor{title}{fg=blue!50!black}
\setbeamercolor{footline}{fg=black!60}
\setbeamertemplate{footline}{%
  \leavevmode\hbox{%
    \begin{beamercolorbox}[wd=.85\paperwidth,ht=2.5ex,dp=1ex,left,leftskip=1em]{footline}\footnotesize\footerlabel\end{beamercolorbox}%
    \begin{beamercolorbox}[wd=.15\paperwidth,ht=2.5ex,dp=1ex,right,rightskip=1em]{footline}\footnotesize\insertframenumber\end{beamercolorbox}}%
  \vskip0pt}
\setbeamertemplate{itemize items}[circle]

\title{Algorithms}
\subtitle{Lecture 2: Binary Search}
\author{Demo Lecturer}
\date{Semester 2, 2026}

\begin{document}

\begin{frame}
  \titlepage
\end{frame}

\begin{frame}{Last Lecture}
  \begin{itemize}
    \item Cost is measured by counting the \textbf{basic operation}
    \item $t(n) \in O(g(n))$ iff $t(n) \le c \cdot g(n)$ for all $n > n_0$
    \item Linear search: worst case $\Theta(n)$ comparisons
  \end{itemize}
  \medskip
  Today: can we do better if the array is \alert{sorted}?
\end{frame}

\begin{frame}{The Problem: Searching a Sorted Array}
  \begin{description}
    \item[Input] an array $A[0..n-1]$ sorted in increasing order, and a key $x$
    \item[Output] an index $j$ with $A[j] = x$, or $-1$ if $x$ is not in $A$
  \end{description}
  \bigskip
  Linear search ignores the order and may scan everything.\\
  The order should let us \emph{rule out many positions with one comparison}.
\end{frame}

\begin{frame}{The Idea: Halve the Search Space}
  Guess a number between 1 and 100; after each guess you are told ``higher'' or ``lower''.
  \medskip
  \begin{center}
  \begin{tabular}{cc}
    \hline
    Guesses so far & Candidates left (at most) \\
    \hline
    0 & 100 \\
    1 & 50 \\
    2 & 25 \\
    3 & 12 \\
    $k$ & $\approx 100 / 2^k$ \\
    \hline
  \end{tabular}
  \end{center}
\end{frame}

\begin{frame}[fragile]{Binary Search: Pseudocode}
\begin{verbatim}
function BinSearch(A, x, n)
    lo <- 0
    hi <- n - 1
    while lo <= hi
        mid <- floor((lo + hi) / 2)
        if A[mid] = x
            return mid
        else if A[mid] < x
            lo <- mid + 1
        else
            hi <- mid - 1
    return -1
\end{verbatim}
\end{frame}

\begin{frame}{Trace: Search for 23}
  $A = [2, 5, 8, 12, 16, 23, 38, 56, 72, 91]$, \quad $x = 23$
  \medskip
  \begin{center}
  \begin{tabular}{cccccl}
    \hline
    Round & lo & hi & mid & $A[\text{mid}]$ & Action \\
    \hline
    1 & 0 & 9 & 4 & 16 & $16 < 23$, so $\text{lo} \leftarrow 5$ \\
    2 & 5 & 9 & 7 & 56 & $56 > 23$, so $\text{hi} \leftarrow 6$ \\
    3 & 5 & 6 & 5 & 23 & found, return 5 \\
    \hline
  \end{tabular}
  \end{center}
\end{frame}

\begin{frame}{Why Must $A$ Be Sorted?}
  The loop maintains an \textbf{invariant}:
  \begin{block}{Invariant}
    If $x$ occurs in $A$, then it occurs in $A[\text{lo}..\text{hi}]$.
  \end{block}
  \begin{itemize}
    \item $A[\text{mid}] < x$: sortedness means everything at or left of mid is $< x$, so discarding it is safe
    \item On an unsorted array the same discard can throw away $x$
  \end{itemize}
\end{frame}

\begin{frame}{Analysis}
  Basic operation: the comparison $A[\text{mid}]$ against $x$.
  \medskip
  \begin{itemize}
    \item \textbf{Best case:} $x$ is at the first midpoint, so 1 comparison
    \item \textbf{Worst case:} the range halves every round; let $C(n)$ be the comparisons:
  \end{itemize}
  \[
    C(n) = C(\lfloor n/2 \rfloor) + 1, \quad C(1) = 1
    \quad\Longrightarrow\quad C(n) = \lfloor \log_2 n \rfloor + 1 \in \Theta(\log n)
  \]
\end{frame}

\begin{frame}{Summary and Homework}
  \begin{itemize}
    \item Binary search needs \textbf{sorted} input and runs in $\Theta(\log n)$ in the worst case
    \item Common bugs: wrong loop condition ($<$ vs $\le$), and forgetting the $\pm 1$ when updating lo / hi
  \end{itemize}
  \bigskip
  \textbf{Homework.} Explain why binary search fails on an unsorted array, with a
  concrete array of length 5 where it returns $-1$ although $x$ is present.
\end{frame}

\end{document}
